% =====================================================================
% limitations.tex — DRAFTED at rb-048 (RB-048), 2026-05-30.
%
% Second bookend; companion of sections/abstract.tex (rb-047, RB-047).
% Mirrors the abstract's voice (distributional / graded / conditional
% by default).  Collects the explicit scope deferrals that every body
% section has already flagged into a single navigable section.
%
% Strength check (mission §3, CLAIM_LEDGER.md): every limitation
% below is a statement about what the live verdict ledger licenses
% and what it does NOT license at the rebuilt strength ceiling.  No
% claim is downgraded or upgraded by this section; the section names
% the rebuild's scope so a reader knows the boundary of every
% headline number.
% =====================================================================

\section{Scope and limitations}
\label{sec:limitations}

This section catalogues the scope of every claim the rebuilt
manuscript makes.  It is organised by lever rather than by section:
each subsection states what the live verdict ledger licenses, what
it does not, and which queued increments would close each remaining
gap.  Nothing here downgrades a rebuilt-strength claim; the purpose
is to make the boundary explicit so a reader can locate the precise
conditional under which any headline number applies.

% ---------------------------------------------------------------
\subsection{The conservation rule is a parameter, not a derivation}
\label{sec:limitations-a3}

The rebuilt model treats the benefit/cost conservation rule as a
power-mean parameter $p$
(Equation~\eqref{eq:conservation-family}, Equation~\eqref{eq:beta-gamma-of-p},
Appendix~\ref{sec:appendix-deriv-a3} Equation~\eqref{eq:a3d-power-mean})
rather than a fixed identity, and reports every headline number at
$p = 1$ (the inherited additive form $\benefit + \cost = 2$) with a
band over the family at $p \in \{0, 0.5, 1\}$
(Section~\ref{sec:extensions-a3}).  Two scoped boundaries follow.

First, the conservation order itself is empirically unconstrained:
nothing in the rebuild fixes which $p$ corresponds to a biological
resource constraint.  The additive choice is the inherited model's;
the multiplicative choice ($p = 0$, $\benefit \cdot \cost = 1$) is
the alternative the original paper noted in its §5.5 but did not
run.  The rebuild does \emph{not} claim either is privileged.  What
it claims is that the central-tendency reframe of $\CF$
(Section~\ref{sec:results-c1}) is robust to the choice — the
per-cell $\Delta \CF = \CF(p = 0) - \CF(p = 1) \le 0$ at every cell
of the $4{,}410$-cell sweep with $0$ reverse flips
(Theorem~\ref{thm:delta-cf-monotone}), so the headline median moves
by $\le 0.001$ but the tail under $\CF < 0.5$ doubles.  Reporting
both branches in the same band is the honest treatment.

Second, the empirical monotonicity $\Delta \CF \le 0$ is licensed by
the rb-016 sim but is \emph{not} a closed-form theorem.  The d'-channel
chain rule of
Appendix~\ref{sec:appendix-deriv-a3}
(Equation~\eqref{eq:a3d-chain-rule},
Equation~\eqref{eq:a3d-dprime-of-p}) yields the analytic statement
$\partial \Rpthree / \partial p \equiv 0$ at $\alpha = 1/\Nloc$
(Proposition~\ref{prop:a3d-rp3-invariance}) — the P3 contribution to
$\CF$ is conservation-form-invariant at the symmetric attention
point — but the joint sign of $\partial \Rpone/\partial p$ and
$\partial \Rpfour/\partial p$ is established only by the envelope
theorem and the empirical 4,410-cell witness, not by a uniform
inequality.  Promoting the chain rule to a closed-form theorem of
the form $|\partial_p R_{\mathrm{P4}}| \le |\partial_p R_{\mathrm{P1}}|$
is an open conjecture.  The joint $(p, \corr > 0)$ band is similarly
deferred: every $p$ sweep above runs at $\corr = 0$, and every
$\corr$ sweep runs at $p = 1$.

% ---------------------------------------------------------------
\subsection{Heterogeneity is bounded, not abolished}
\label{sec:limitations-a2-a8}

The rebuild lifts the inherited model's single-$\Rsens$ scalar to a
per-location vector $\boldsymbol{r}$
(Equation~\eqref{eq:d-prime-hetero}) and the inherited
1-dimensional uncued allocation to a full $\Nloc$-dimensional
simplex
(Section~\ref{sec:extensions-a8}, the policy-space lift of
\texttt{er\_full\_policy}), and shows that both extensions leave the
inherited model's optima essentially intact under the additive
conservation rule.  Two boundaries remain.

First, the between-preparation reading of $\Rsens$
(Section~\ref{sec:model}) is the one in which the model statement is
binding; the within-display reading
(\cite{mcadams_maunsell1999_v4_tuning,reynolds_heeger2009_normalization,carrasco2011_visual_attention_25y})
is exercised empirically by the $\Rsens_i$ sweep
(Section~\ref{sec:extensions-a2}) but is not promoted to a closed-form
theorem.  The empirical statement is that the $C_2$ peak coordinates
and the contested-$\CF$ corner are invariant to within $\le 1{\times}10^{-5}$
across linear spreads $\le 30\%$
(Table~\ref{tab:a2-rb021-summary}), so the rebuild's headline
numbers transport.  A closed-form perturbation expansion
$\Delta R = O(\mathrm{Var}(\boldsymbol{r}))$ around the homogeneous
optimum is queued as a future derivation increment.

Second, the new A8 conditional — the $\Nloc$-dimensional uncued
policy binds in a benefit-dominant symmetric-stress corner under
multiplicative conservation
(Section~\ref{sec:extensions-a8}, Finding F2: $\Delta R = +2.79\times
10^{-3}$ at $\corr = 0$, amplified to $+3.68\times 10^{-3}$ at
$\corr = 0.2$, at $\valid = 1/\Nloc$, $\val = 1$, $\Rsens = 10$,
variant~A) — is exhibited at a single $\Rsens$, a single $(\valid,
\val)$ cell, a single $\Nloc$, and variant~A only.  The closed-form
A8-binding-onset locus in $(\Rsens, p, \valid)$ is not derived; a
Slepian-style monotonicity argument on the joint no-FA integrand
under simplex perturbation is the natural substrate but has not been
attempted.  The variant-B replication of the F2 cell, the finer
$\Rsens$ grid around $\Rsens \in \{5, 10, 20\}$, and the
$\Nloc > 4$ generalisation are queued as separate sims.

% ---------------------------------------------------------------
\subsection{The decision-noise lever is deferred, not denied}
\label{sec:limitations-a6}

The rebuilt three-lever decomposition
(Definition~\ref{def:three-levers}: criterion, sensitivity,
decorrelation) does \emph{not} include attention-coupled decision
noise as a fourth lever, even though the reviewer's continuing
analysis of the homogeneous-decision-rule premise (A6 in the live
ledger) flagged it as a candidate third channel that the criterion
fraction would mis-book to ``attention''.  At the current live label
(\texttt{WEAKLY-SUPPORTED}), the A6 attack-vector evidence does not
yet shift any rebuilt headline number — fixed per-location decision
noise is benign (absorbed into effective $\dprime$), and the
attention-coupled-noise crack is suggestive rather than decisive.
The rebuild therefore holds the lever inventory at three, with the
explicit caveat that an attention-coupled-noise channel
$\sigma_i(\alpha)$ would, if landed, enter alongside $\corr$ as a
fourth axis of the lever decomposition and would shift the
interpretation of $\CF$ in the cost-dominant regime.  The model and
sim increments that would wire this lever are queued and gated on
the A6 verdict moving past \texttt{WEAKLY-SUPPORTED}; the rebuild
will fold the lever in by symmetry with the $\corr$ channel when
the verdict licenses it.

% ---------------------------------------------------------------
\subsection{Assumptions retained as explicit scope}
\label{sec:limitations-retained}

Three assumptions of the inherited model are carried into the
rebuilt model unchanged, because the live verdict ledger has not
attacked them and they were not blocking any rebuilt-strength claim.
Each is named here so the boundary is explicit.

\emph{A4 (no learning dynamics).}  The observer is taken to have
already discovered the optimal policy; the rebuilt manuscript is a
normative statement about the policy itself, not about how it is
acquired.  Learning speed and reliability — and any further bias
they introduce toward the simpler criterion-based lever — are
outside the rebuild's scope.  Folding a learning dynamic into the
$(\corr, p, \boldsymbol{r}, \text{policy-space})$ extensions of the
rebuilt model is a substantial separate program.

\emph{A5 (transfer-function family).}  The four monotonic forms
$h \in \{a, \sqrt{a}, a^{0.3}, a^{2}\}$ used by the inherited paper
exhaust the qualitative landscape probed by the rebuilt sims; the
rebuild's headline cell fixes $h = \sqrt{\,\cdot\,}$ for
consistency with the inherited $C_2$ peak figure.  Other monotonic
families (sigmoidal, threshold, contrast-response-style) are
plausible biological extensions and are not tested.  The rebuilt
model's mechanism statements (Definition~\ref{def:three-levers},
Equation~\eqref{eq:r-dagger}, Proposition~\ref{prop:r-dagger-invariance})
are stated in terms of generic $f$ and do not depend on the
particular family; the empirical bands of
Sections~\ref{sec:results-c1}--\ref{sec:results-c4} do.

\emph{A7 (reward variants).}  The two reward conventions
(variant~A: $\CR = \valid \cdot \val + (1 - \valid)$;
variant~B: $\CR = 1$) bracket the natural extremes the inherited
paper studied.  The rebuilt manuscript reports the headline numbers
under variant~A throughout, with variant~B explicitly cited as a
sensitivity (Table~\ref{tab:cf-distribution},
Table~\ref{tab:cf-quadrants}, the variant-B $\CF$ tail at $0.30$,
the variant-B equal-split criticality of
Table~\ref{tab:a2-rb021-summary}).  Other reward structures
(continuous-valued, asymmetric cost) are not tested.  The
variant-axis caveats below
(Section~\ref{sec:limitations-followups}) belong here in spirit; we
list them with the other queued replications for compactness.

% ---------------------------------------------------------------
\subsection{Equicorrelation, $\corr$-envelope, and queued follow-ups}
\label{sec:limitations-followups}

Three further boundaries collect the remaining sims and derivations
the rebuild has queued but not run.

\emph{Equicorrelation and the $\corr$-envelope.}  The A1 channel
(Section~\ref{sec:model-rho-channel},
Equation~\eqref{eq:equicorr-cov},
Equation~\eqref{eq:pnofa-rho}) assumes a single scalar $\corr$
shared by every cross-location pair.  Structured covariances —
distance-dependent decay, feature-similarity-dependent grouping,
inter-areal cross-correlations
\cite{RuffCohen2016,Srinath2021} — break the exact one-factor
Gauss--Hermite reduction and are out of scope.  The $\corr$ grid is
tested at $\corr \in \{0, 0.1, 0.2, 0.3, 0.4\}$ with the central
empirical anchor $\corr = 0.2$ from
Cohen--Maunsell~\cite{CohenMaunsell2009}; the $\corr \ge 0.5$ regime
is plausible only under heavy attention-mediated decorrelation and
is not exercised.  The closed-form $\rdagger(\val; \corr)$ drift
of Appendix~\ref{sec:appendix-deriv-c2}
(Proposition~\ref{prop:r-dagger-rho}, Table~\ref{tab:r-dagger-rho-drift})
reports the $\corr = 0 \to 0.2$ single step; a finer $\corr$-grid
characterisation of the drift rate $d\rdagger/d\corr$ as a function
of $\val$ is queued.

\emph{Variant-B caveats.}  The rebuild reports headline numbers
under variant~A throughout; the variant-B replications of the
$\corr$-channel $\CF$ flatness probe
(Section~\ref{sec:model-upper-bound}), the $C_2$ peak under
$\corr > 0$, the $C_3$ iso-$\VDA$ contour band, the anti-cue $C_4$
inversion (Section~\ref{sec:results-c4}), the A2 cost-dominant
$\Delta R$ suppression
(Section~\ref{sec:extensions-a2}, Finding 2), and the new A8 F2
binding (Section~\ref{sec:extensions-a8}) are queued as separate
sims.  The rb-002 finding that variant-B $\CF$ is essentially flat
in $\corr$ at the headline cell makes the variant axis genuinely
informative as a sensitivity, not redundant; the rebuild flags this
explicitly rather than treating variant~B as a footnote.

\emph{Grid sharpening.}  Several quantitative thresholds are pinned
to the grid step of their underlying sim and are queued for
sharpening if a referee asks: the $C_3$ high-$\valid$ threshold
$\valid^{*}$
(Section~\ref{sec:results-c3}; pinned to $\pm 0.05$, finer grid
queued); the $C_4$ anti-cue onset $\valid^{*}$ at $\val = 1$,
$\Nloc = 4$ (Section~\ref{sec:results-c4}; pinned to $\pm 0.05$,
finer grid queued); the A8 binding onset in $\Rsens$ near
$\{5, 10, 20\}$ (Section~\ref{sec:extensions-a8}; pinned to a
single $\Rsens = 10$ probe, finer grid queued); the closed-form
$\CF < 0.5$ boundary in $(\valid, \Rsens)$
(Section~\ref{sec:results-c1}; conjectured diagonal in
$(\log \Rsens, \valid)$, derivation queued); the dormant-cell
amplification band at $\Rsens = 0.3$, $(\valid \approx 0.7,
\val = 10)$ (Section~\ref{sec:results-c3}, where $\VDA = 0.0007$ at
$\corr = 0$ lifts to $0.0676$ at $\corr = 0.2$ — a $96\times$
amplification that is the most striking single qualitative finding
of the $C_3$ sweep, queued for a clean closeup); and the
Slepian-style analytic locus for the cell-wise $\partial \VDA /
\partial \corr$ sign-flip surface
(Section~\ref{sec:model-upper-bound}, Table~\ref{tab:a1cw-summary};
cell-wise crossover $\Rsens^{\times} \approx 0.79$ pinned
empirically, formal proposition queued).  None of these queued
items would change a rebuilt-strength headline; each would either
tighten a boundary or promote an empirical observation to an
analytic statement.

% ---------------------------------------------------------------
\subsection{What the rebuild does not claim}
\label{sec:limitations-not-claimed}

The rebuilt manuscript makes no claim about the neural
implementation of any of the three levers — neither the criterion
shift, nor the sensitivity re-allocation, nor the decorrelation
channel is mapped onto a particular cortical circuit, recurrent
dynamics, learning rule, or architectural primitive.  It makes no
claim about the empirical $\corr$ values that human or animal
observers exhibit under cued change detection beyond citing the
Cohen--Maunsell~\cite{CohenMaunsell2009} $r_{SC} \approx 0.2$ V4
anchor as the central operating point of the $\corr$ envelope.  It
does not predict any particular biological observer's deviation
from the normative policy.  And it does not address attention
dynamics over time (within-trial buildup, sequential effects across
trials, learning trajectories) — every statement above is about a
stationary observer at a fixed parameter cell.  The rebuild's
positive contribution is mathematical: a more honestly bounded
description of when value-directed attention matters, with the
boundary made explicit lever by lever, so that any downstream
empirical or computational program inherits a clearly scoped
substrate rather than a single uniformly-stated headline.
